\documentclass[11pt]{letter}
\usepackage[margin=1in]{geometry}
\usepackage{parskip}
\usepackage[hidelinks]{hyperref}

\signature{Rishit Puri}
\address{Rishit Puri\\
ORCID 0009-0004-6215-9997\\
\texttt{rishitpuri15@gmail.com}}

\begin{document}

\begin{letter}{Editor-in-Chief\\
IEEE Journal of Biomedical and Health Informatics}

\opening{Dear Editor,}

Please consider the enclosed manuscript, \emph{``What Polyp Segmentation
Benchmarks Cannot Measure: Small Lesions, Domain Shift, and a Noise Floor
Larger Than Most Reported Gains,''} for publication as a regular paper in the
IEEE Journal of Biomedical and Health Informatics.

Automated polyp segmentation is now almost universally benchmarked as a
single mean Dice score on one train/test split of Kvasir-SEG, and reported
values cluster within a few points of one another. Rather than proposing
another architecture, this manuscript asks what that protocol is capable of
measuring. We train five architectures spanning four backbone families under
an identical recipe, evaluate them with five-fold cross-validation over all
1000 Kvasir-SEG images and zero-shot on four external datasets, and measure
the protocol's own variance by repeating identical configurations.

Three findings seem to us of direct interest to the journal's readership:

\begin{enumerate}
\item \textbf{Clinically small lesions are missed about half the time.} For
lesions occupying under 1\% of the frame, per-lesion detection falls to
0.44--0.59 from 0.86--0.89 for larger lesions. Of the 123 such images in the
external sets, 44 (36\%) are detected by \emph{no} architecture, and 16 elicit
an empty mask from every model. Kvasir-SEG contains six such images out of
1000, so the field's primary benchmark is close to blind to the failure mode
that dominates clinical misses.

\item \textbf{The evaluation protocol carries a noise floor comparable to
reported gains.} Repeating identical training configurations yields a
test-Dice standard deviation of 0.0053, and 4 of 10 pairwise in-domain
architecture comparisons fall below it. We trace the variance not to random
seeds or non-deterministic CUDA kernels---neither of which removes it---but to
checkpoint selection on a flat validation curve, where identical runs select
epochs 11.3 apart on average and differences in validation Dice carry no
information about differences in test Dice ($r=0.004$).

\item \textbf{Test-set choice determines resolving power.} Expressed in units
of that noise floor, the same five architectures span 3.9$\sigma$ in domain
but 21.6$\sigma$ on ETIS-Larib, so cross-dataset evaluation distinguishes
methods that in-domain evaluation cannot.
\end{enumerate}

We also report, as a worked example of the first problem, a subgroup finding
from our own single-split analysis---an apparent robustness of U-Net++ to
flat lesions---that did not survive cross-validation. We include it because
it is precisely the error the conventional protocol invites.

The work is intended to be immediately actionable for practitioners:
it recommends reporting a measured noise floor, including datasets that
contain genuinely small lesions, reporting per-lesion detection alongside
region overlap, and avoiding single-checkpoint model selection on small
validation sets. All code, cross-validation splits and per-image results for
all 88 training runs are publicly released
(\url{https://github.com/rishitpuri/polyp-segmentation-evaluation},
archived at \url{https://doi.org/10.5281/zenodo.22841360}).

I confirm that this manuscript is original, has not been published
previously, and is not under consideration by any other journal. There are no
prior or concurrent submissions of overlapping content. The study used only
publicly available, de-identified datasets, so institutional ethics approval
was not required. I am employed as a software engineer by Karl Storz
Endoscopy; the company had no role in the conception, design, execution,
analysis or reporting of this work, provided no funding or resources for it,
and the study uses exclusively public datasets and open-source software. This
is declared in the manuscript's competing interests statement. I am the sole
author and have no other competing interests to declare.

% Optional: ScholarOne will ask for suggested reviewers. Nominate three or
% four researchers who publish on polyp segmentation benchmarking or on
% evaluation methodology in medical imaging, and with whom you have no
% conflict (no shared institution, supervision or recent co-authorship).
% Do not nominate anyone you have not actually read.

Thank you for your time and consideration.

\closing{Sincerely,}

\end{letter}
\end{document}
